\documentclass[12pt,a4paper]{article}
\usepackage{goyabase}

\usepackage{longtable}
\usepackage{calc}

% Pandoc compatibility
\providecommand{\tightlist}{%
  \setlength{\itemsep}{0pt}\setlength{\parskip}{0pt}}
\newcommand{\passthrough}[1]{#1}

\title{\textbf{Cassandra en Docker - UT8 NoSQL}}
\author{Departamento de Informática}
\date{\today}

\usepackage[spanish, es-noshorthands, es-tabla]{babel}
\usepackage[autostyle]{csquotes}
\begin{document}

\maketitle

\section{Cassandra en Docker - UT8 NoSQL}\label{cassandra-en-docker---ut8-nosql}

Este entorno permite trabajar con \textbf{Apache Cassandra}, una base de datos \textbf{Orientada a Columnas} diseñada para gestionar volúmenes masivos de datos con escrituras ultra-rápidas.

\subsection{El Ecosistema de Datos de GloboTour}\label{ecosistema}
Para este laboratorio, simularemos tres áreas críticas de la plataforma GloboTour que requieren la potencia de Cassandra. A continuación se muestran ejemplos reales de la información almacenada:

\begin{enumerate}
    \item \textbf{Módulo de Inteligencia de Mercado (\texttt{busquedas\_por\_usuario})}:
    \begin{itemize}
        \item \textit{Dato Real:} El usuario \texttt{user\_7} buscó el destino \enquote{Londres} el \texttt{2026-04-15} desde la IP \texttt{192.168.1.54}.
    \end{itemize}
    \item \textbf{Módulo de Telemetría IoT (\texttt{sensor\_readings})}: 
    \begin{itemize}
        \item \textit{Dato Real:} El sensor \texttt{SENSOR-SALA-A} del hotel \texttt{H-MAD-01} registró una temperatura de \texttt{22.5 ºC} a las \texttt{10:00 AM}.
    \end{itemize}
    \item \textbf{Módulo de Marketing Dinámico (\texttt{offer\_stats} y \texttt{flash\_offers})}: 
    \begin{itemize}
        \item \textit{Dato Real:} La oferta \texttt{OFERTA-PARIS-2024} ha recibido \texttt{120} visualizaciones. Existe una promoción flash \texttt{FLASH-01} con un \texttt{25\%} de descuento que caducará automáticamente tras el tiempo definido por su TTL.
    \end{itemize}
\end{enumerate}

\subsection{Componentes}\label{componentes}

\begin{enumerate}
    \item \textbf{Cassandra Server}: Motor de base de datos en el puerto \texttt{9042}.
    \item \textbf{CloudBeaver}: Interfaz web profesional para gestión de BBDD en el puerto \texttt{8978}.
\end{enumerate}

\subsection{Instrucciones de uso}\label{instrucciones-de-uso}

\subsubsection{Configuración del Entorno (Docker Compose)}\label{configuracion-del-entorno}

Este archivo define los servicios necesarios para la práctica:

\begin{lstlisting}[language=YAML]
services:
  cassandra:
    image: cassandra:latest
    container_name: cassandra-server
    ports:
      - "9042:9042"
    volumes:
      - cassandra_data:/var/lib/cassandra
      - ./src:/src
    restart: always

  cloudbeaver:
    image: dbeaver/cloudbeaver:latest
    container_name: cloudbeaver
    ports:
      - "8978:8978"
    volumes:
      - cloudbeaver_data:/opt/cloudbeaver/workspace
    depends_on:
      - cassandra
    restart: always

volumes:
  cassandra_data:
  cloudbeaver_data:
\end{lstlisting}

\subsubsection{Puesta en marcha}\label{puesta-en-marcha}

Para arrancar el servicio:

\begin{lstlisting}[language=bash]
docker compose up -d
\end{lstlisting}

\emph{Nota: Cassandra puede tardar hasta 60 segundos en responder.}

Para importar los datos y el esquema inicial:

\begin{lstlisting}[language=bash]
docker exec -i cassandra-server cqlsh < src/datos_iniciales.cql
\end{lstlisting}

Para acceder a la interfaz gráfica:
\begin{itemize}
    \item URL: \url{http://localhost:8978}
    \item \textbf{Configuración inicial}: Sigue el asistente y crea un usuario admin.
    \item \textbf{Conexión}: Host: \texttt{cassandra-server} | Puerto: \texttt{9042} | Keyspace: \texttt{globotour}.
\end{itemize}


Para abrir la consola:

\begin{lstlisting}[language=bash]
docker exec -it cassandra-server cqlsh
\end{lstlisting}

\end{document}
